% This file is included by uet_thesis.tex.
% Replace blue placeholders and remove example text before submission.

\chapter{Design Specification}
\label{chap:design}
% =============================================================================

\section{System Context, Architecture, and Interfaces \applicability{Development, Hybrid, Game; Research Software Prototype}}
The system boundary contains the Flutter web/Android client, the ASP.NET Core
marketplace API, the PostgreSQL/PostGIS database, Redis, SignalR, and the
asynchronous Python/FastAPI AI service. Customers and Contractors submit and
discover lots, communicate, negotiate, and reserve them. Admins manage users,
categories, thresholds, flagged listings, and AI overrides. AWS S3 stores listing
photographs. The AI service returns classification, defect, condition,
consistency, and price signals to the decision engine.

\begin{table}[H]
\centering
\caption{Major interfaces}
\label{tab:major-interfaces}
\begin{tabularx}{\textwidth}{@{}p{0.16\textwidth}p{0.22\textwidth}X X@{}}
\toprule
\textbf{ID} & \textbf{Interface} & \textbf{Purpose and exchanged information} & \textbf{Protocol or format} \\
\midrule
INT-01 & Flutter client--ASP.NET Core API & Authentication, listings, search, reservations, and Admin operations & REST/JSON \\
INT-02 & ASP.NET Core--FastAPI AI service & Listing photos and metadata exchanged for AI assessment results & Asynchronous API \\
INT-03 & ASP.NET Core--SignalR & Real-time chat and negotiation messages & SignalR \\
INT-04 & Platform--PostgreSQL/PostGIS and Redis & Persistent marketplace data, geospatial queries, cache, and presence & Database/cache protocols \\
INT-05 & Platform--AWS S3 & Listing-photo upload and retrieval & Object storage API \\
\bottomrule
\end{tabularx}
\end{table}

\section{Research Methodology and Experimental Design \applicability{Research, Hybrid}}
\subsection{Study or Experimental Design}
The project will evaluate the integrated system using held-out image data and
Admin review decisions. Each submitted lot passes through the same classification,
defect, condition, consistency, and price stages. Results are compared with
validated labels, sourced reference prices, and human review outcomes.

\subsection{Population, Sample, Participants, or Dataset}
The AI pipeline will use existing public building-defect datasets, including
YOLO-format datasets, SDNET2018, and MBDD2025, supplemented by a small
team-collected validation set of marketplace-style phone photographs. New-material
prices will be collected from published Pakistani construction-pricing sources.
Surplus depreciation will be informed by structured interviews with 5--10 Lahore
dealers and contractors. User requirements will be validated through interviews
with 4--6 contractors/suppliers and a brief questionnaire.

\subsection{Variables and Data Collection Techniques}

\begin{table}[H]
\centering
\caption{Variables and data collection}
\label{tab:variables}
\begin{tabularx}{\textwidth}{@{}p{0.16\textwidth}p{0.23\textwidth}X X@{}}
\toprule
\textbf{ID} & \textbf{Variable or construct} & \textbf{Operational definition or measure} & \textbf{Collection source or instrument} \\
\midrule
VAR-01 & \placeholder{Variable or construct} & \placeholder{How it is defined or measured} & \placeholder{Instrument, sensor, dataset field, interview, etc.} \\
VAR-02 & Material classification accuracy & Correct predicted category across Bricks, Steel/Rebar, Tiles/Marble, and Wood & Held-out image set \\
VAR-03 & Defect detection performance & Precision and recall for each configured defect class & Labelled public datasets and real-photo validation set \\
VAR-04 & Price-estimate error & MAE and RMSE against reference prices & Published Pakistani price sources \\
VAR-05 & Decision agreement & Agreement between automated publish/flag result and Admin decision & Admin review outcomes \\
\bottomrule
\end{tabularx}
\end{table}

\subsection{Analysis and Evaluation Protocol}
Evaluation will report material-classification accuracy and confusion patterns,
defect-detection precision and recall per defect class, agreement between the
automated publish/flag decision and Admin review decisions, and price-estimate
MAE/RMSE against sourced reference prices. The real-photo validation set will be
kept separate from training data.



\section{Behavioural and Interaction Design \applicability{Development, Hybrid, Game; Research if applicable}}
\begin{guidance}
Include only the models needed to explain the project's behaviour and interaction.
Each UML diagram must use standard UML notation~\citep{omguml251}, remain readable, and be followed by a
brief explanation of important flows, alternatives, or design decisions.
\end{guidance}

\subsection{Use-Case Model}
The principal use cases are Create Lot Listing, Search Listings, Chat and
Negotiate, Reserve Lot, Review Flagged Listing, Override AI Decision, and Ask the
AI Listing Assistant.

\begin{table}[H]
\centering
\caption{Use-case specification}
\label{tab:use-case-specification}
\begin{tabularx}{\textwidth}{@{}p{0.24\textwidth}X@{}}
\toprule
\textbf{Field} & \textbf{Description} \\
\midrule
Use-case ID and name & UC-01: Create and verify lot listing \\
Primary actor & Contractor/Seller \\
Preconditions & Authenticated account and lot details with 3--5 photographs \\
Trigger & Seller submits a listing \\
Main flow & Store lot; run AI pipeline; publish if checks pass or queue for Admin review \\
Alternative or exception flows & Low confidence, disagreement, or Mixed condition routes to Admin \\
Postconditions & Listing is published, rejected, or awaiting review \\
Related requirements & FR-01--FR-05, NFR-01--NFR-02 \\
\bottomrule
\end{tabularx}
\end{table}

\subsection{Activity Model}
\begin{guidance}
Model an important workflow using actions, control flows, decisions with guards,
parallel paths where required, and clear initial and final nodes. Use swimlanes when
responsibility changes are important.
\end{guidance}
\diagramplaceholder{Activity diagram}{Example workflow: submit request, check availability, approve or reject, and notify the student.}

\subsection{State-Machine Model \applicability{State-dependent systems and games}}
\begin{guidance}
Model the lifecycle of one important stateful entity. Show states, triggering
events, guards where needed, actions, and initial and final states.
\end{guidance}
\diagramplaceholder{State-machine diagram}{Example entity: Appointment; states may include Requested, Approved, Rescheduled, Completed, and Cancelled.}

\subsection{Sequence Model}
\begin{guidance}
Show how an important scenario is realized through time-ordered messages among
actors, system components, and external services. Use activation bars and combined
fragments such as \texttt{alt}, \texttt{opt}, or \texttt{loop} where appropriate.
\end{guidance}
\diagramplaceholder{Sequence diagram}{Example scenario: a student submits an appointment request and receives confirmation.}

\section{Structural Design \applicability{Development, Hybrid, Game; Research if applicable}}
\begin{guidance}
Describe the static structure needed to realize the requirements. Keep models at a
level appropriate to the project and avoid including generated framework classes
that do not communicate a design decision.
\end{guidance}

\subsection{Class or Domain Model}
\begin{guidance}
Show important domain concepts or software classes, relevant attributes and
operations, relationships, multiplicities, and inheritance where justified.
\end{guidance}
\diagramplaceholder{Class or domain diagram}{Example classes: Student, Adviser, Appointment, AvailabilitySlot, and Notification.}

\subsection{Component Model}
\begin{guidance}
Show deployable or replaceable software components, their provided and required
interfaces, and dependencies.
\end{guidance}
\diagramplaceholder{Component diagram}{Example components: Web Client, Appointment Service, Authentication Service, and Data Store.}

\subsection{Package Model \applicability{If required}}
\begin{guidance}
Use a package diagram when the organization and dependencies of modules or layers
need clarification. Omit it when the component or class model already communicates
the structure adequately.
\end{guidance}
\diagramplaceholder{Package diagram}{Example packages: presentation, application, domain, and infrastructure.}

\subsection{Deployment Model \applicability{Deployed, distributed, networked, or multi-device systems}}
\begin{guidance}
Show execution nodes, devices, environments, deployed artifacts, and communication
paths. Distinguish physical or virtual nodes from software components.
\end{guidance}
\diagramplaceholder{Deployment diagram}{Example nodes: user device, application server, database server, and external notification service.}

\section{Data Design and Management}
\label{sec:data-design}

\subsection{Database, Schema, and Persistence Design \applicability{Development, Hybrid, Game, if applicable}}
\begin{guidance}
Describe persistent entities, identifiers, relationships, integrity constraints,
and important indexing or retention decisions. Use an ER diagram, schema, or other
notation appropriate to the selected database model.
\end{guidance}
\diagramplaceholder{Data model or schema}{Example entities: User, AvailabilitySlot, Appointment, and NotificationLog.}

\subsection{Research Data Management \applicability{Research, Hybrid}}
\begin{guidance}
Explain how research data will be collected or acquired, named, stored, protected,
anonymized or pseudonymized, versioned, retained, shared, and disposed of. Identify
access controls and responsible team members.
\end{guidance}

\begin{table}[H]
\centering
\caption{Research data management plan}
\label{tab:data-management}
\begin{tabularx}{\textwidth}{@{}p{0.22\textwidth}X X@{}}
\toprule
\textbf{Data category} & \textbf{Management method} & \textbf{Protection, retention, or access rule} \\
\midrule
\placeholder{Raw/processed data, code, consent records, model artifacts} & \placeholder{Format, location, naming, and versioning} & \placeholder{Access, retention, sharing, and disposal} \\
\bottomrule
\end{tabularx}
\end{table}

\section{Algorithm or Model Design \applicability{Research, Hybrid; Development or Game if a core algorithm or AI model is used}}
\begin{table}[H]
\centering
\caption{Core algorithm or model}
\label{tab:algorithm-model}
\begin{tabularx}{\textwidth}{@{}p{0.18\textwidth}X@{}}
\toprule
\textbf{Item} & \textbf{Description} \\
\midrule
Name and purpose & CLIP, YOLO, condition engine, consistency checker, and XGBoost price model for listing verification. \\
Inputs and outputs & Representative listing photographs, description, quantity, location, and reference price; material, defects, condition, consistency result, and estimated surplus price. \\
Processing steps & Classify material; detect visible defects; aggregate photo results into a lot grade; compare description and image signals; estimate price; publish or flag using decision thresholds. \\
Assumptions and limits & Assessment covers representative photos and only defect classes present and validated in available datasets. \\
Related requirements & FR-02--FR-04, NFR-02 \\
\bottomrule
\end{tabularx}
\end{table}

\section{User Interface and Accessibility Design \applicability{User-facing Development, Hybrid, and Game projects; Research prototypes if applicable}}
\begin{guidance}
Present the navigation structure and key wireframes or prototypes. Explain how the
design supports usability, accessibility, responsive layout, error prevention,
feedback, and consistency. Do not include every minor screen.
\end{guidance}
\diagramplaceholder{Navigation structure or user flow}{Example: Sign In, Dashboard, Available Slots, Appointment Details, and Notifications.}
\diagramplaceholder{Key wireframe or prototype}{Annotate important controls, content, validation, feedback, and accessibility decisions.}

\section{Security, Privacy, and Safety Design \applicability{If applicable}}
\begin{guidance}
Explain the design controls used to satisfy the corresponding requirements in
Chapter~\ref{chap:requirements}. Address relevant threats, authentication,
authorization, data protection, logging, secure failure, privacy, and safety. Do
not make general claims such as ``the system is secure'' without specific controls.
\end{guidance}

\begin{table}[H]
\centering
\caption{Security, privacy, and safety controls}
\label{tab:security-controls}
\begin{tabularx}{\textwidth}{@{}p{0.16\textwidth}X X p{0.19\textwidth}@{}}
\toprule
\textbf{Design ID} & \textbf{Risk or requirement} & \textbf{Design control} & \textbf{Related requirement(s)} \\
\midrule
SEC-01 & Unauthorized access to appointment records & Role-based authorization is checked at the service layer before records are returned or modified. & FR-XX, NFR-XX \\
SEC-02 & \placeholder{Threat, privacy concern, or safety hazard} & \placeholder{Specific design control} & \placeholder{Requirement IDs} \\
\bottomrule
\end{tabularx}
\end{table}

\section{Design Decisions and Rationale}
\begin{guidance}
Record significant choices, alternatives considered, and reasons for selection.
Focus on decisions that materially affect quality, cost, risk, maintainability, or
project feasibility. Do not repeat the tools list from Section~\ref{sec:tools}.
\end{guidance}

\begin{table}[H]
\centering
\caption{Design decisions}
\label{tab:design-decisions}
\begin{tabularx}{\textwidth}{@{}p{0.14\textwidth}X X X@{}}
\toprule
\textbf{ID} & \textbf{Decision} & \textbf{Alternatives considered} & \textbf{Rationale and trade-off} \\
\midrule
DD-01 & \placeholder{Selected design decision} & \placeholder{Feasible alternatives} & \placeholder{Evidence-based reason and accepted limitation} \\
\bottomrule
\end{tabularx}
\end{table}

% =============================================================================
